\chapter{Mathematical Background}
\label{app:math}

This appendix provides a quick reference for the mathematical concepts used throughout the book.

\section{Number Theory}

\subsection{Modular Arithmetic}

For integers $a, b, n$ with $n > 0$:
\[
a \equiv b \pmod{n} \iff n \mid (a - b)
\]

Properties:
\begin{align*}
a \equiv b \pmod{n},\; c \equiv d \pmod{n} &\implies a+c \equiv b+d \pmod{n} \\
&\implies ac \equiv bd \pmod{n} \\
&\implies a^k \equiv b^k \pmod{n}
\end{align*}

\subsection{Euler's Totient Function}

\[
\varphi(n) = n \prod_{p \mid n} \left(1 - \frac{1}{p}\right)
\]

Properties:
\begin{itemize}
    \item $\varphi(p) = p-1$ for prime $p$
    \item $\varphi(p^k) = p^k - p^{k-1}$
    \item $\varphi(mn) = \varphi(m)\varphi(n)$ if $\gcd(m,n)=1$
    \item Euler's theorem: $a^{\varphi(n)} \equiv 1 \pmod{n}$ for $\gcd(a,n)=1$
\end{itemize}

\subsection{Chinese Remainder Theorem}

If $n_1, \ldots, n_k$ are pairwise coprime, then the system
\[
x \equiv a_i \pmod{n_i}, \quad i=1,\ldots,k
\]
has a unique solution modulo $N = n_1 n_2 \cdots n_k$.

Gauss's constructive solution:
\[
x = \sum_{i=1}^k a_i N_i M_i \pmod{N}, \quad
N_i = \frac{N}{n_i},\quad M_i \equiv N_i^{-1} \pmod{n_i}
\]

\subsection{Quadratic Residues and Legendre Symbol}

For odd prime $p$ and $a \not\equiv 0 \pmod{p}$:
\[
\left(\frac{a}{p}\right) = 
\begin{cases}
1 & \text{if } a \text{ is a quadratic residue mod } p \\
-1 & \text{if } a \text{ is a quadratic non-residue mod } p
\end{cases}
\]

Euler's criterion: $\left(\frac{a}{p}\right) \equiv a^{(p-1)/2} \pmod{p}$

\textbf{Law of Quadratic Reciprocity} (Gauss's "Golden Theorem"):
For distinct odd primes $p, q$:
\[
\left(\frac{p}{q}\right)\left(\frac{q}{p}\right) = (-1)^{\frac{p-1}{2}\frac{q-1}{2}}
\]

\textbf{Supplements}:
\begin{align*}
\left(\frac{-1}{p}\right) &= (-1)^{(p-1)/2} = \begin{cases} 1 & p \equiv 1 \pmod{4} \\ -1 & p \equiv 3 \pmod{4} \end{cases} \\
\left(\frac{2}{p}\right) &= (-1)^{(p^2-1)/8} = \begin{cases} 1 & p \equiv \pm 1 \pmod{8} \\ -1 & p \equiv \pm 3 \pmod{8} \end{cases}
\end{align*}

\subsection{Gauss Sums}

Quadratic Gauss sum:
\[
g(p) = \sum_{a=0}^{p-1} e^{2\pi i a^2/p} = 
\begin{cases}
\sqrt{p} & p \equiv 1 \pmod{4} \\
i\sqrt{p} & p \equiv 3 \pmod{4}
\end{cases}
\]

General Gauss sum with Dirichlet character $\chi$:
\[
\tau(\chi) = \sum_{a=0}^{n-1} \chi(a) e^{2\pi i a/n}
\]

\subsection{Binary Quadratic Forms}

A binary quadratic form: $f(x,y) = ax^2 + bxy + cy^2$ with discriminant $\Delta = b^2 - 4ac$.

Gauss's composition law makes the set of equivalence classes of primitive forms of discriminant $\Delta < 0$ into an abelian group (the class group).

\section{Analysis}

\subsection{Gaussian Integral}

\[
\int_{-\infty}^{\infty} e^{-x^2} dx = \sqrt{\pi}
\]

Generalized:
\[
\int_{-\infty}^{\infty} e^{-ax^2+bx+c} dx = \sqrt{\frac{\pi}{a}} e^{c + b^2/(4a)}, \quad \Re(a) > 0
\]

\subsection{Error Function}

\[
\erf(x) = \frac{2}{\sqrt{\pi}} \int_0^x e^{-t^2} dt, \quad
\erfc(x) = 1 - \erf(x) = \frac{2}{\sqrt{\pi}} \int_x^{\infty} e^{-t^2} dt
\]

Properties:
\begin{itemize}
    \item $\erf(-x) = -\erf(x)$
    \item $\erf(0) = 0$, $\lim_{x\to\infty} \erf(x) = 1$
    \item $\erf'(x) = \frac{2}{\sqrt{\pi}} e^{-x^2}$
    \item Series: $\erf(x) = \frac{2}{\sqrt{\pi}} \sum_{n=0}^{\infty} \frac{(-1)^n x^{2n+1}}{n!(2n+1)}$
    \item Asymptotic: $\erfc(x) \sim \frac{e^{-x^2}}{x\sqrt{\pi}}\left(1 - \frac{1}{2x^2} + \cdots\right)$
\end{itemize}

\subsection{Normal Distribution}

Standard normal PDF and CDF:
\[
\phi(x) = \frac{1}{\sqrt{2\pi}} e^{-x^2/2}, \quad
\Phi(x) = \frac{1}{2}\left(1 + \erf\left(\frac{x}{\sqrt{2}}\right)\right)
\]

General $\mathcal{N}(\mu, \sigma^2)$:
\[
f(x) = \frac{1}{\sigma\sqrt{2\pi}} e^{-(x-\mu)^2/(2\sigma^2)}, \quad
F(x) = \Phi\left(\frac{x-\mu}{\sigma}\right)
\]

\subsection{Arithmetic-Geometric Mean (AGM)}

For $a_0, b_0 > 0$:
\[
a_{n+1} = \frac{a_n + b_n}{2}, \quad b_{n+1} = \sqrt{a_n b_n}
\]
Converge quadratically to common limit $M(a_0, b_0)$.

Gauss's identity:
\[
\frac{1}{M(1, k')} = \frac{2}{\pi} K(k), \quad k' = \sqrt{1-k^2}
\]
where $K(k)$ is the complete elliptic integral of the first kind.

Gauss's constant: $G = 1/M(1, \sqrt{2}) \approx 0.8346268$

\subsection{Brent-Salamin Algorithm for $\pi$}

$a_0 = 1,\; b_0 = 1/\sqrt{2},\; s_0 = 1/2$

For $n \ge 0$:
\begin{align*}
a_{n+1} &= (a_n + b_n)/2 \\
b_{n+1} &= \sqrt{a_n b_n} \\
s_{n+1} &= s_n - 2^n (a_n - b_n)^2
\end{align*}

Then $\pi = \dfrac{4 a_N^2}{1 - s_N}$ with quadratic convergence.

\subsection{Theta Functions}

Jacobi theta functions ($q = e^{\pi i \tau}$):
\begin{align*}
\vartheta_1(z|\tau) &= -i \sum_{n=-\infty}^{\infty} (-1)^n q^{(n+1/2)^2} e^{2\pi i (n+1/2)z} \\
\vartheta_2(z|\tau) &= \sum_{n=-\infty}^{\infty} q^{(n+1/2)^2} e^{2\pi i (n+1/2)z} \\
\vartheta_3(z|\tau) &= \sum_{n=-\infty}^{\infty} q^{n^2} e^{2\pi i n z} \\
\vartheta_4(z|\tau) &= \sum_{n=-\infty}^{\infty} (-1)^n q^{n^2} e^{2\pi i n z}
\end{align*}

Jacobi's triple product identity:
\[
\vartheta_1(z|\tau) = 2 q^{1/4} \sin(\pi z) \prod_{n=1}^{\infty} (1 - q^{2n})(1 - 2 q^{2n} \cos(2\pi z) + q^{4n})
\]

Modular transformation:
\[
\vartheta_3(0|-1/\tau) = \sqrt{-i\tau}\, \vartheta_3(0|\tau)
\]

\subsection{Dedekind Eta Function}

\[
\eta(\tau) = q^{1/24} \prod_{n=1}^{\infty} (1 - q^n), \quad q = e^{2\pi i \tau}
\]

Modular transformation:
\[
\eta(-1/\tau) = \sqrt{-i\tau}\, \eta(\tau)
\]

\section{Linear Algebra}

\subsection{Gaussian Elimination}

Forward elimination with partial pivoting, then back substitution.
Complexity: $O(n^3)$.

\subsection{LU Decomposition}

$PA = LU$ where $P$ is permutation, $L$ unit lower triangular, $U$ upper triangular.

\subsection{Cholesky Decomposition}

For symmetric positive definite $A$: $A = LL^T$ with $L$ lower triangular.

\subsection{QR Decomposition}

$A = QR$ with $Q$ orthogonal, $R$ upper triangular.
Householder reflections: $H = I - 2uu^T$.

\subsection{Condition Number}

$\kappa(A) = \|A\| \|A^{-1}\|$ for a given norm.
For 2-norm: $\kappa_2(A) = \sigma_{\max}/\sigma_{\min}$.

\subsection{Least Squares}

Minimize $\|Ax - b\|_2^2$.
Normal equations: $A^T A x = A^T b$.
QR solution: $Rx = Q^T b$ (more stable).

\subsection{Gaussian Quadrature}

For weight function $w(x)$ on $[a,b]$, nodes $x_i$ are zeros of orthogonal polynomial $p_n(x)$.
Weights $w_i = \int_a^b \frac{p_n(x)}{(x-x_i)p_n'(x_i)} w(x) dx$.
Exact for polynomials of degree $\le 2n-1$.

Golub-Welsch algorithm: eigenvalues of Jacobi matrix give nodes; first eigenvector components squared give weights.

\section{Statistics}

\subsection{Error Propagation (Delta Method)}

For $y = f(x_1, \ldots, x_n)$ with uncertainties $\sigma_i$:
\[
\sigma_y^2 \approx \sum_{i=1}^n \left(\frac{\partial f}{\partial x_i}\right)^2 \sigma_i^2
\]

For correlated variables with covariance matrix $\Sigma_x$:
\[
\Sigma_y = J \Sigma_x J^T
\]
where $J$ is the Jacobian of $f$.

\subsection{Gaussian Process}

A Gaussian process is a collection of random variables $\{f(x) : x \in \mathcal{X}\}$ such that any finite subset has a multivariate normal distribution.

Specified by mean function $m(x)$ and covariance kernel $k(x, x')$:
\[
f \sim \mathcal{GP}(m, k)
\]

GP regression predictive distribution:
\begin{align*}
\mu_* &= m(X_*) + K(X_*, X) [K(X, X) + \sigma_n^2 I]^{-1} (y - m(X)) \\
\Sigma_* &= K(X_*, X_*) - K(X_*, X) [K(X, X) + \sigma_n^2 I]^{-1} K(X, X_*)
\end{align*}

Log marginal likelihood:
\[
\log p(y|X, \theta) = -\frac{1}{2} y^T K_\theta^{-1} y - \frac{1}{2} \log |K_\theta| - \frac{n}{2} \log 2\pi
\]

Common kernels:
\begin{itemize}
    \item RBF: $k(x,x') = \sigma^2 \exp(-\|x-x'\|^2/(2\ell^2))$
    \item Matérn $\nu=3/2$: $k(x,x') = \sigma^2 (1 + \frac{\sqrt{3}r}{\ell}) e^{-\sqrt{3}r/\ell}$
    \item Periodic: $k(x,x') = \sigma^2 \exp(-\frac{2\sin^2(\pi|x-x'|/p)}{\ell^2})$
\end{itemize}

\section{Geometry}

\subsection{First Fundamental Form}

For surface $\mathbf{r}(u,v)$:
\[
E = \mathbf{r}_u \cdot \mathbf{r}_u, \quad
F = \mathbf{r}_u \cdot \mathbf{r}_v, \quad
G = \mathbf{r}_v \cdot \mathbf{r}_v
\]

\subsection{Second Fundamental Form}

Unit normal $\mathbf{N} = \frac{\mathbf{r}_u \times \mathbf{r}_v}{\|\mathbf{r}_u \times \mathbf{r}_v\|}$:
\[
L = \mathbf{r}_{uu} \cdot \mathbf{N}, \quad
M = \mathbf{r}_{uv} \cdot \mathbf{N}, \quad
N = \mathbf{r}_{vv} \cdot \mathbf{N}
\]

\subsection{Curvature}

Gaussian curvature: $K = \frac{LN - M^2}{EG - F^2}$

Mean curvature: $H = \frac{EN - 2FM + GL}{2(EG - F^2)}$

\textbf{Theorema Egregium}: $K$ depends only on $E, F, G$ and their derivatives (intrinsic).

Christoffel symbols:
\begin{align*}
\Gamma_{11}^1 &= \frac{GE_u - 2FF_u + FE_v}{2(EG-F^2)} \\
\Gamma_{12}^1 &= \frac{GE_v - FG_u}{2(EG-F^2)} \\
\Gamma_{22}^1 &= \frac{2GF_v - GG_u - FG_v}{2(EG-F^2)} \\
\Gamma_{11}^2 &= \frac{2EF_u - EE_v - FE_u}{2(EG-F^2)} \\
\Gamma_{12}^2 &= \frac{EG_u - FE_v}{2(EG-F^2)} \\
\Gamma_{22}^2 &= \frac{EG_v - 2FF_v + FG_u}{2(EG-F^2)}
\end{align*}

\subsection{Gauss-Bonnet Theorem}

For compact surface $S$ with boundary $\partial S$:
\[
\int_S K\, dA + \int_{\partial S} \kappa_g\, ds = 2\pi \chi(S)
\]

Discrete version for triangle mesh:
\[
\sum_i (2\pi - \theta_i) = 2\pi \chi
\]
where $\theta_i$ is sum of angles at vertex $i$.

\section{Astronomy}

\subsection{Orbital Elements}

Six classical elements:
\begin{itemize}
    \item $a$ -- semi-major axis
    \item $e$ -- eccentricity
    \item $i$ -- inclination
    \item $\Omega$ -- longitude of ascending node
    \item $\omega$ -- argument of periapsis
    \item $M_0$ -- mean anomaly at epoch
\end{itemize}

\subsection{Kepler's Equation}

$M = E - e \sin E$, solved iteratively:
\[
E_{n+1} = E_n - \frac{E_n - e\sin E_n - M}{1 - e\cos E_n}
\]

\subsection{Two-Body Problem}

$\ddot{\mathbf{r}} = -\frac{\mu}{r^3} \mathbf{r}$, $\mu = G(M+m)$.

Conserved quantities:
\begin{itemize}
    \item Angular momentum: $\mathbf{h} = \mathbf{r} \times \dot{\mathbf{r}}$
    \item Energy: $\mathcal{E} = \frac{1}{2}\|\dot{\mathbf{r}}\|^2 - \frac{\mu}{r}$
    \item Laplace-Runge-Lenz vector: $\mathbf{e} = \frac{1}{\mu}\dot{\mathbf{r}} \times \mathbf{h} - \frac{\mathbf{r}}{r}$
\end{itemize}

\subsection{Universal Variables (f and g functions)}

$\mathbf{r} = f \mathbf{r}_0 + g \dot{\mathbf{r}}_0$,
$\dot{\mathbf{r}} = \dot{f} \mathbf{r}_0 + \dot{g} \dot{\mathbf{r}}_0$

With universal anomaly $\chi$:
\begin{align*}
f &= 1 - \frac{\chi^2}{r_0} C(z) \\
g &= \Delta t - \frac{\chi^3}{\sqrt{\mu}} S(z) \\
\dot{f} &= \frac{\sqrt{\mu}}{r r_0} \chi (z S(z) - 1) \\
\dot{g} &= 1 - \frac{\chi^2}{r} C(z)
\end{align*}
where $z = \alpha \chi^2$, $C(z)$ and $S(z)$ are Stumpff functions.